\emph{(i)} Write $C^{k} = \operatorname{Hom}(\Lambda^{k}\mathfrak{s}, \mathfrak{s})$ with the Chevalley-Eilenberg differential $(dc)(x_{0}, \dots, x_{k}) = \sum_{a}(-1)^{a}[x_{a}, c(\dots\hat{x}_{a}\dots)] + \sum_{a<b}(-1)^{a+b}c([x_{a}, x_{b}], \dots\hat{x}_{a}\dots\hat{x}_{b}\dots)$, where $\mathfrak{s} = \mathfrak{sl}(4,\C) \cong \mathfrak{sl}(\CS)$ by Proposition \claimref{correspondence.spin-representation}. The Cartan subalgebra $\mathfrak{h}$ acts on $C^{k}$ by the Lie derivative $L_{H}$, which commutes with $d$ and satisfies $L_{H} = d\,\iota_{H} + \iota_{H}\,d$ (Cartan's formula, with $(\iota_{H}c)(x_{1}, \dots) = c(H, x_{1}, \dots)$). On the subcomplex of torus weight $\lambda \neq 0$ choose $H$ with $\lambda(H) \neq 0$; for a cocycle $c$ there, $c = \lambda(H)^{-1}L_{H}c = d(\lambda(H)^{-1}\iota_{H}c)$, so that subcomplex is acyclic and $H^{k}(\mathfrak{s}, \mathfrak{s})$ is the cohomology of the weight-zero subcomplex, the cochains mapping each tuple of basis vectors into the weight space of the sum of their weights. In the basis $E_{ij}$ ($i \neq j$), $H_{k} = E_{kk} - E_{k+1,k+1}$, the weight-zero cochain spaces in degrees $0, 1, 2, 3$ have dimensions $3, 21, 87, 273$ and the check computes the ranks of $d_{0}, d_{1}, d_{2}$ exactly as $3, 18, 69$, with $d_{2}d_{1} = 0$; so $\dim H^{1} = 21 - 18 - 3 = 0$ and $\dim H^{2} = 87 - 69 - 18 = 0$. (This is Whitehead's second lemma for $\mathfrak{sl}(4)$, recomputed.) Rigidity follows by the standard order-by-order argument (Gerstenhaber's deformation theory, cited from memory for attribution): if $\mu \equiv [\,,\,]$ modulo $\hbar^{m}$ after a redefinition, the Jacobi identity for $\mu$ at order $\hbar^{m}$ says that $\mu_{m}$ is a 2-cocycle, hence $\mu_{m} = d\phi_{m}$, and conjugating by $1 - \hbar^{m}\phi_{m}$ (with the sign fixed by $d$) removes it modulo $\hbar^{m+1}$; the product of these maps is $\Phi$. Then with $Q = q\circ\Phi^{-1}$, $Q(X)\star Q(Y) - Q(Y)\star Q(X) = \hbar\,q(\mu(\Phi^{-1}X, \Phi^{-1}Y)) = \hbar\,q(\Phi^{-1}[X, Y]) = \hbar\,Q([X,Y])$, the undeformed relation of Proposition \claimref{correspondence.symmetrised-quadratics-palatial-algebra}(2). The ordering tensor of Lemma \claimref{correspondence.ordering-filtration}(2) is, under (CQ), $\mu$ read through $q$; so its class modulo redefinitions of the generators is zero.

\emph{(ii)} By Lemma \claimref{correspondence.chiral-planes-summand-spaces}(1) one of $\mathfrak{g}^{\pm}$ acts faithfully on $\Som$ and trivially on $\Spi$, the other the reverse, so in a basis of $\CS$ adapted to $\Som \oplus \Spi$ they are the two diagonal $\mathfrak{sl}(2)$ blocks, and $[\mathfrak{g}^{+}, \mathfrak{g}^{-}] = 0$ as in Proposition \claimref{correspondence.chirality-split-ordering-insensitivity-asd-sector}(1); which block is called self-dual does not matter below. With $(h_{+}, e_{+}, f_{+})$ the standard basis of $\mathfrak{g}^{+}$, the map $\phi(f_{+}) = e_{+}$, $\phi = 0$ on $h_{+}, e_{+}$, on $\mathfrak{g}^{-}$ and on the rest of $\mathfrak{s}$, satisfies $[\operatorname{ad}h_{+}, \phi] = 4\phi$ on $\mathfrak{g}^{+}$ and $\operatorname{tr}(\phi(x)y) = \operatorname{tr}(x\phi(y))$ on $\mathfrak{g}^{+}$: it is the highest-weight vector of the spin-two summand of $\operatorname{End}\mathfrak{g}^{+} \cong \mathfrak{g}^{+}\otimes\mathfrak{g}^{+} = \mathbf{1}\oplus\mathbf{3}\oplus\mathbf{5}$ (the symmetric traceless part; complete reducibility as in Fact \claimref{correspondence.imported-complete-reducibility} is not needed beyond this explicit vector). For $M, N \in \mathfrak{g}^{-}$, $\phi M = \phi N = \phi[M,N] = 0$, so $d\phi(M,N) = 0$; for $M \in \mathfrak{g}^{-}$, $N \in \mathfrak{g}^{+}$, $d\phi(M,N) = [M, \phi N] = 0$ since $\phi N \in \mathfrak{g}^{+}$ commutes with $M$, and $[N, \phi M] = \phi[M,N] = 0$. On $\mathfrak{g}^{+}$, $d\phi(h_{+}, f_{+}) = [h_{+}, e_{+}] - [f_{+}, 0] - \phi(-2f_{+}) = 4e_{+} \neq 0$, and all values on $(\mathfrak{g}^{+}, \mathfrak{g}^{+})$ lie in $\mathfrak{g}^{+}$. For $M \in \mathfrak{g}^{-}$ and $Y \in \mathfrak{p}$, $d\phi(M, Y) = [M, \phi Y] - [Y, \phi M] - \phi[M, Y] = 0$, because $\phi Y = \phi M = 0$ and $[\mathfrak{g}^{-}, \mathfrak{p}] \subset \mathfrak{p}$, on which $\phi$ vanishes. For $N \in \mathfrak{g}^{+}$ and $Y \in \mathfrak{p}$, $d\phi(N, Y) = -[Y, \phi N] - \phi[N, Y]$, and for $Y, Y' \in \mathfrak{p}$, $d\phi(Y, Y') = -\phi[Y, Y']$, nonzero whenever $[Y, Y']$ has an $f_{+}$ component; the check exhibits nonzero values on both. So $\phi$ is supported on $\mathfrak{g}^{+}$ but $d\phi$ is not: the redefinition changes brackets with the non-Lorentz generators too, which the anti-self-dual consistency check does not see. Under $\mu = [\,,\,] + \hbar\,d\phi$, the defects of words in $q(\mathfrak{g}^{-})$ are those of Proposition \claimref{correspondence.chirality-split-ordering-insensitivity-asd-sector}(2)(a), and the googly slot still commutes with $q(\mathfrak{g}^{-})$ to first order, so the consistency check is passed; the check verifies all of this exactly, together with $(1 + t\phi)([M,N] + t\,d\phi(M,N)) - [(1 + t\phi)M, (1 + t\phi)N] = t^{2}(\phi\,d\phi(M,N) - [\phi M, \phi N])$ at $t = \tfrac17$, which exhibits $d\phi$ as the first-order transport of the bracket by $1 + \hbar\phi$. A spin-two tensor of the self-dual factor has the five components of a self-dual Weyl spinor; that this is the right identification is a reading, not proved here.

\emph{(iii)} By (i), under (CQ) a deformation defined over an open set of the spacetime (or of the twistor data) is, on each small patch $U$ where the order-by-order construction runs, the redefinition $\Phi_{U}$. Invariant content therefore resides only in the transition data $\Phi_{U}^{-1}\Phi_{V}$, which is argued, not proved, to be where part (b) would have to locate the googly data, by analogy with the nonlinear graviton, whose anti-self-dual data is the patching of a locally standard twistor space (Fact \claimref{correspondence.imported-nonlinear-graviton} records the construction; it is not used in the proof). Deformations in which $q(M)\star q(N) - q(N)\star q(M)$ leaves the span of the quadratics, as a deformation of the full palatial algebra generally may, are outside (CQ) and untouched by this argument.
